%% CV General - Álvaro López Álvarez (Español)
%% CV completo, no orientado a una oferta concreta. Extensión flexible.
%% Compilar con: lualatex main_general_es.tex
%% (pdflatex suele fallar en instalaciones modernas de MiKTeX con errores de fontawesome5.)

\documentclass[11pt,a4paper,sans]{moderncv}
\moderncvstyle{banking}
\moderncvcolor{blue}

% Fuerza el nombre, apellidos y los títulos de sección al azul de moderncv
% (color1). Por defecto, el estilo banking en lualatex+MiKTeX los deja en
% negro, lo cual desentona con el resto de la paleta azul.
\renewcommand*{\firstnamestyle}[1]{{\fontsize{34}{36}\bfseries\upshape\color{color1}#1}}
\renewcommand*{\lastnamestyle}[1]{{\fontsize{34}{36}\bfseries\upshape\color{color1}#1}}
\renewcommand*{\sectionstyle}[1]{{\sectionfont\color{color1}#1}}

\usepackage[utf8]{inputenc}
\usepackage{needspace}
\usepackage{hyperref}
\hypersetup{
    colorlinks=true,
    linkcolor=blue,
    filecolor=magenta,
    urlcolor=blue,
    pdftitle={Álvaro López Álvarez - CV},
    pdfpagemode=FullScreen,
}
\usepackage[scale=0.80]{geometry}
\usepackage{import}
\usepackage{eso-pic}

% datos personales
\name{Álvaro}{López Álvarez}
\address{Valencia, España}{}{}
\phone[mobile]{+34 640764553}
\email{alvaro.lopez.dev@gmail.com}
\extrainfo{\href{https://www.linkedin.com/in/alvarolopez-dev}{LinkedIn}, \href{https://github.com/alerone}{GitHub}, \href{https://alerone.github.io/portfolio}{Portfolio}}

% El estilo banking de moderncv (variante de cabecera 3) no conecta \photo -
% se ignora silenciosamente. Se superpone la foto directamente, anclada a la
% esquina de la página para que sea independiente de la cabecera centrada.
\AddToShipoutPictureFG{%
  \ifnum\value{page}=1
    \AtPageUpperLeft{%
      \hspace*{\dimexpr\paperwidth-2.3cm-1.3cm\relax}%
      \raisebox{\dimexpr-1.3cm-2.3cm\relax}{\fboxsep0pt\fboxrule0.6pt\fcolorbox{color1}{white}{\includegraphics[width=2.3cm]{photo}}}%
    }%
  \fi
}

\begin{document}

\makecvtitle

% ============================================================
%     PERFIL PROFESIONAL
% ============================================================

\vspace{6pt}
\small{Ingeniero de Software con una sólida formación en Ingeniería Informática y experiencia práctica en desarrollo backend, sistemas distribuidos, ingeniería de producto full-stack y herramientas de IA agéntica basadas en LLM. Lenguajes principales: Golang y TypeScript, con conocimientos prácticos de Kotlin, Java, Python y SQL. Cómodo asumiendo funcionalidades de extremo a extremo, desde arquitecturas backend orientadas a eventos hasta frontends en React, con especial interés en ingeniería de sistemas de bajo nivel -- compiladores, intérpretes y protocolos de red construidos desde cero. Máster en Ingeniería Informática finalizado (nota media 8,8/10) y nivel B2 de inglés certificado, mientras desarrolla proyectos personales de forma activa.}

% ============================================================
%     COMPETENCIAS CLAVE
% ============================================================

\section{Competencias Clave}
\vspace{1pt}
\begin{itemize}

\item \textbf{Backend y Lenguajes}: Golang, TypeScript/JavaScript, Java, Python, Kotlin, SQL, GDScript, C

\item \textbf{Sistemas Distribuidos}: Docker, NATS, Apache Kafka, Apache APISIX, microservicios orientados a eventos

\item \textbf{Desarrollo Full-Stack}: React, PostgreSQL, APIs REST, OAuth 2.0/JWT

\item \textbf{Ingeniería de IA/LLM}: LangChain, LangGraph, API de Gemini, API de OpenAI, pipelines RAG

\item \textbf{Herramientas y Prácticas}: Git, GitHub Actions, Docker Compose, Testcontainers, TDD, Firebase, Supabase, Vercel, Render, Android Studio, NeoVim

\end{itemize}

% ============================================================
%     EXPERIENCIA PROFESIONAL
% ============================================================

\section{Experiencia Profesional}
\vspace{3pt}
\begin{itemize}

% --- Puesto más reciente ---
\needspace{5\baselineskip}
\item{\cventry{Feb 2026--Jun 2026}{Desarrollador I+D}{Instituto Tecnológico de Informática (ITI)}{Valencia, España}{}{\vspace{1pt}
\begin{itemize}
    \item Implementación de mejoras basadas en LLM para un motor de búsqueda
    \item Desarrollo de un agente multi-paso con LangGraph que convierte preguntas en lenguaje natural en consultas SQL validadas de solo lectura, para resolver políticas de acceso a datos sobre un esquema de gobernanza de datos en PostgreSQL
    \item Diseño agéntico orientado a la seguridad: selección de esquema/tablas, generación estructurada de SQL, validación de ejecución, reintento ante errores
    \item Trabajo dentro de una aplicación centrada en la gestión, gobernanza y soberanía de datos industriales
\end{itemize}}}

\vspace{3pt}

% --- Puesto anterior ---
\item{\cventry{Feb 2024--Jun 2024}{Desarrollador Full-Stack (Prácticas)}{Instituto Tecnológico de Informática (ITI)}{Valencia, España}{}{\vspace{1pt}
\begin{itemize}
    \item Desarrollo y mantenimiento de la aplicación corporativa de gestión de personal (Intranet)
    \item Implementación de nuevas funcionalidades de frontend con React y TypeScript, mejorando la usabilidad interna
    \item Contribución al desarrollo backend, garantizando una integración fluida entre servicios
\end{itemize}}}

\end{itemize}

% ============================================================
%     PROYECTOS PERSONALES DESTACADOS
% ============================================================

\section{Proyectos Personales Destacados}
\vspace{1pt}
\begin{itemize}

\item \textbf{AIDucate} -- \textit{Go, Gin, PostgreSQL, React, TypeScript, Tailwind, API de Gemini}: Plataforma SaaS impulsada por IA que genera planes de estudio, contenido y ejercicios personalizados. Backend en Go/Gin con autenticación JWT y persistencia en PostgreSQL; frontend en React/TypeScript desplegado en Vercel, backend en Render. \textit{(En desarrollo)}

\item \textbf{ValenMove} -- \textit{Kotlin, MVVM, JSoup}: Aplicación Android nativa que unifica el transporte público de Valencia (EMT, Metrovalencia, Valenbisi) en una sola interfaz; ingeniería inversa de APIs públicas combinada con web scraping y el portal de datos abiertos del ayuntamiento para llegadas en tiempo real. Trabajo de Fin de Grado, calificación 9/10.

\item \textbf{FaaS con Golang} -- \textit{Go, Docker, NATS JetStream, Apache APISIX}: Plataforma distribuida de Function-as-a-Service en la que workers en Go ejecutan funciones arbitrarias empaquetadas en imágenes Docker. Colas de trabajo con NATS JetStream para una distribución fiable de tareas, gateway con Apache APISIX para autenticación JWT y limitación de tasa, y escalado horizontal mediante réplicas de Docker Compose.

\item \textbf{Plataforma de Venta de Entradas con Kafka y Spring Boot} -- \textit{Java, Spring Boot 3, Apache Kafka, Testcontainers}: Microservicios orientados a eventos (reserva, pago, notificación, auditoría) con patrón outbox, colas de mensajes fallidos (dead-letter), procesamiento idempotente y pruebas de integración con Testcontainers.

\item \textbf{Monkey Interpreter} -- \textit{Go}: Intérprete de árbol de sintaxis para un lenguaje propio, construido a partir del libro \textit{Writing an Interpreter in Go}: lexer, parser Pratt, AST, evaluador y REPL escritos desde cero.

\end{itemize}

% ============================================================
%     FORMACIÓN ACADÉMICA
% ============================================================

\section{Formación Académica}
\vspace{1pt}
\begin{itemize}

\needspace{5\baselineskip}
\item{\cventry{2024--2026}{Máster Universitario en Ingeniería Informática (MUIINF)}{Universitat Politècnica de València (UPV)}{Valencia, España}{}{\vspace{1pt}
Finalizado (septiembre de 2026), nota media 8,8/10. Temario: sistemas distribuidos, computación de altas prestaciones, ciberseguridad, sistemas empotrados; asignaturas cursadas: Sistemas Empotrados y Ubicuos, Redes y Seguridad.
}}

\vspace{3pt}

\needspace{5\baselineskip}
\item{\cventry{2020--2024}{Grado en Ingeniería Informática, itinerario de Ingeniería del Software}{Universitat Politècnica de València (UPV)}{Valencia, España}{}{\vspace{1pt}
Nota media: 8,88/10. Trabajo de Fin de Grado (TFG): ``ValenMove -- Aplicación Unificada de Transporte Público'' (calificación 9/10).
}}

\end{itemize}

% ============================================================
%     IDIOMAS
% ============================================================

\section{Idiomas}
\vspace{1pt}
\begin{itemize}
\item Español (nativo), catalán (nativo), inglés (B2 certificado, conversación fluida).
\end{itemize}

% ============================================================
%     REFERENCIAS
% ============================================================

\section{Referencias}
\vspace{1pt}
\begin{itemize}
\item{Disponibles a petición.}
\end{itemize}

\end{document}
